\documentclass[12pt]{article}
\usepackage[margin=1in]{geometry}

\usepackage{amsmath,amssymb,amsfonts}
\usepackage{booktabs}
\usepackage{graphicx}
\usepackage{hyperref}
\usepackage{xcolor}
\usepackage{multirow}
\usepackage{subcaption}

\newcommand{\bx}{\mathbf{x}}
\newcommand{\bX}{\mathbf{X}}
\newcommand{\bdx}{\dot{\mathbf{x}}}
\newcommand{\bXi}{\boldsymbol{\Xi}}
\newcommand{\bTheta}{\boldsymbol{\Theta}}
\newcommand{\RR}{\mathbb{R}}
\newcommand{\norm}[1]{\left\lVert #1 \right\rVert}
\newcommand{\abs}[1]{\left\lvert #1 \right\rvert}
\newcommand{\CLW}{\text{CLW}}

\graphicspath{{outputs/figures/}}

\begin{document}

\title{Sparse Equation Discovery on the Chen--Lin--White Tokamak Model: \\
A Controlled Benchmark of Noise, Derivative Estimation, and Library Design}

\author{Lucas Perrier}
\date{}

\maketitle

\begin{abstract}
We study sparse equation discovery on the Chen--Lin--White (CLW) reduced tokamak turbulence model as a controlled benchmark rather than as a new algorithmic contribution.
The CLW system is a four-dimensional chaotic ODE with a rational nonlinearity, making it a useful test case for assessing how standard SINDy behaves on a physically motivated but still tractable problem.
Using STLSQ with a physics-informed candidate library, we first confirm an idealized upper bound: with clean states and exact right-hand-side derivatives, the true eight-term model is recovered to machine precision.
We then evaluate three derivative regimes---analytic right-hand-side evaluation on noisy states, raw finite differences, and Savitzky--Golay smoothed finite differences---across six state-noise levels and five random seeds.
Recovery is consistently strong for $\eta \le 10^{-2}$ and deteriorates rapidly near $\eta = 10^{-1}$, with smoothed finite differences giving the best low-noise coefficient accuracy but not extending the usable noise range.
Two library ablations sharpen the interpretation: an extended degree-2 library still recovers the clean system but admits additional false positives under noise, while an incomplete library missing $(PZ/S)\sin C$ produces a persistent coefficient-error floor of about $0.48$ even at negligible noise.
We also show that a simple penalized training-error criterion selects models near the threshold-path knee and behaves similarly to held-out validation, but neither that criterion nor BIC removes a non-monotonic sample-size trend caused by correlated, physically spurious features on the attractor.
Because CLW is chaotic, we argue that support recovery, coefficient error, and short-horizon vector-field agreement are more informative than long-horizon rollout fidelity.
The result is a reproducible baseline for future work on equation discovery in fusion-relevant reduced models.
\end{abstract}

\noindent\textbf{Keywords:} sparse identification of nonlinear dynamics, tokamak turbulence, equation discovery, chaotic systems, SINDy, Chen--Lin--White model

\bigskip

%% ============================================================
\section{Introduction}
\label{sec:introduction}

Understanding when sparse equation discovery succeeds on nonlinear plasma models is more useful, at this stage, than claiming a new discovery pipeline for fusion dynamics.
High-fidelity plasma simulations remain expensive, and low-dimensional reduced models continue to matter for mechanism studies, control-oriented modeling, and methodology development~\cite{Chen1984,Horton1999}.
If sparse regression is to play a role in that workflow, practitioners need realistic baselines on systems that are more structured than textbook examples but still simple enough to admit careful error accounting.

The Sparse Identification of Nonlinear Dynamics (SINDy) framework~\cite{Brunton2016} offers a natural starting point.
Given state data $\bX$ and derivatives $\bdx$, SINDy seeks a sparse coefficient matrix $\bXi$ such that $\bdx \approx \bTheta(\bX)\,\bXi$ for a chosen library $\bTheta(\bX)$ of candidate functions.
Its appeal is interpretability: the output is an explicit set of governing equations rather than a black-box predictor.
Its practical difficulty is equally clear.
Recovery quality depends strongly on how derivatives are estimated, on whether the library contains the right terms, and on how model quality is judged when the underlying dynamics are chaotic.

Those three issues are usually explored separately.
Derivative estimation under noise has been studied extensively~\cite{Schaeffer2017,Rudy2017}; library design and overcomplete dictionaries are a known challenge for sparse regression~\cite{Mangan2016}; and chaotic systems are notoriously hard to evaluate with long-horizon rollouts alone~\cite{Strogatz2015}.
For a plasma-relevant reduced model, however, an honest end-to-end baseline is still useful: not to claim that standard SINDy solves the problem, but to identify the regimes in which it works, the regimes in which it fails, and the failure modes that remain after careful tuning.

We therefore use the Chen--Lin--White (CLW) model~\cite{Chen1984} as a benchmark problem.
CLW couples a resistive tearing mode and an ion temperature gradient mode through four ODEs in the variables $P$, $S$, $Z$, and $C$.
It is chaotic at the parameter setting considered here, yet small enough that the true coefficient matrix is known exactly and the consequences of library misspecification can be measured directly.
The rational term $(PZ/S)\sin C$ is particularly useful for benchmarking because it is physically meaningful, non-polynomial, and difficult to approximate if the library is chosen poorly.

The central claim of this paper is deliberately modest.
We do not propose a new optimizer, a new differentiator, or a new plasma model.
Instead, we provide a reproducible benchmark study of baseline STLSQ-SINDy on CLW, with enough ablations to make the results informative for both nonlinear-dynamics readers and plasma-modeling readers.
The main contributions are:

\begin{enumerate}
\item We establish an idealized upper bound by showing exact recovery from clean states and exact derivatives when the candidate library contains the true terms.
\item We compare three practically relevant derivative regimes under measurement noise and identify a consistent failure threshold near $\eta = 10^{-1}$ for this system.
\item We quantify the separate effects of library overspecification and library misspecification using extended and incomplete libraries.
\item We test whether the main observations are artifacts of threshold tuning by examining threshold paths, held-out validation, sample size, BIC selection, and SFD hyperparameters.
\item We argue that, for chaotic systems such as CLW, equation-level recovery metrics are the primary quantities of interest and that long-horizon trajectory agreement should be treated as supporting evidence only.
\end{enumerate}

The intended scope is therefore narrower, but more defensible, than a general-purpose SINDy paper.
This is a benchmark paper about one physically motivated chaotic system, written to be useful precisely because it is explicit about what is idealized, what is robust, and what is not.


%% ============================================================
\section{Related Work}
\label{sec:related}

\paragraph{Sparse equation discovery.}
The SINDy framework was introduced by Brunton, Proctor, and Kutz~\cite{Brunton2016}, building on compressed sensing and sparse regression.
Extensions address noisy data through integral formulations~\cite{Schaeffer2017}, weak formulations~\cite{Messenger2021}, and ensemble methods~\cite{Fasel2022}.
Rudy et al.~\cite{Rudy2017} extended the approach to partial differential equations.
Champion et al.~\cite{Champion2019} combined SINDy with autoencoders for coordinate discovery in high-dimensional systems.
We deliberately restrict our study to the baseline STLSQ algorithm in order to establish a clear reference point against which more advanced methods can be compared on the CLW system.

\paragraph{Derivative estimation.}
Accurate derivative estimation is a known bottleneck for equation discovery from noisy data.
Total variation regularization~\cite{Chartrand2011}, Savitzky--Golay filtering~\cite{Savitzky1964}, and spectral differentiation~\cite{Schaeffer2017} have been proposed.
The PySINDy package~\cite{deSilva2020,Kaptanoglu2022} provides a unified interface to these methods through its \texttt{SmoothedFiniteDifference} differentiator.
Our contribution is to compare three derivative regimes on the same chaotic system under identical noise and library conditions.

\paragraph{Reduced plasma models.}
The CLW system was derived by Chen, Lin, and White~\cite{Chen1984} as a reduced description of resistive tearing mode--ITG coupling in tokamak geometry.
It has served as a compact model for nonlinear mode interaction and chaotic plasma dynamics~\cite{Horton1999}.
Sparse system identification has also been explored on other plasma-relevant settings~\cite{Dam2017,Kaptanoglu2020}, but a focused baseline study of derivative noise, library choice, and chaos-aware evaluation on CLW has been missing.

\paragraph{Evaluation of discovered equations in chaotic systems.}
The exponential divergence of nearby trajectories in chaotic systems implies that long-horizon forecast error is a poor metric for model quality~\cite{Strogatz2015}.
Short-horizon comparisons and statistical measures (e.g., Lyapunov exponents, attractor statistics) are more appropriate~\cite{Pathak2018,Vlachas2018}.
We build on this insight by normalizing trajectory comparisons by the Lyapunov time and focusing evaluation primarily on coefficient-level metrics.

\paragraph{Benchmark studies.}
Several works have benchmarked SINDy variants on canonical systems (Lorenz, Lotka--Volterra, etc.); see, e.g., the comparison in~\cite{Fasel2022}.
Our study extends this benchmarking tradition to a plasma-physics-motivated system whose rational nonlinearity $(PZ/S)\sin C$ and four-dimensional chaotic attractor present challenges not found in the standard test suite.


%% ============================================================
\section{The Chen--Lin--White System}
\label{sec:system}

The CLW system~\cite{Chen1984} is a four-dimensional autonomous ODE:
\begin{align}
\dot{P} &= P - 2\,Z\,S\,\cos C, \label{eq:clw_P} \\
\dot{S} &= -\Gamma_d\,S + Z\,P\,\cos C, \label{eq:clw_S} \\
\dot{Z} &= -\gamma_z\,Z + 2\,P\,S\,\cos C, \label{eq:clw_Z} \\
\dot{C} &= d - \frac{P\,Z}{S}\,\sin C, \label{eq:clw_C}
\end{align}
where the state variables are:
\begin{itemize}
\item $P$: tearing mode amplitude,
\item $S$: sound-wave amplitude,
\item $Z$: ITG mode amplitude,
\item $C$: relative phase between the modes.
\end{itemize}
The parameters are the sound-wave damping rate $\Gamma_d$, the ITG mode growth rate $\gamma_z$, and the frequency mismatch $d$.
Throughout this work, we fix $\Gamma_d = 2.0$, $\gamma_z = 0.8$, and $d = 2.0$, which places the system in a chaotic regime.

The system has $k^\star = 8$ nonzero coefficients across the four equations: one linear term per equation for $\dot{P}$, $\dot{S}$, and $\dot{Z}$; the three cross-coupling cosine terms; a constant forcing in $\dot{C}$; and the rational sine term in $\dot{C}$.
The rational term $(PZ/S)\sin C$ is noteworthy because it involves division by a state variable, requiring careful library design to ensure representability.

\subsection{Chaotic Dynamics and the Lyapunov Exponent}

At the chosen parameters, the CLW system is chaotic.
We estimate the maximal Lyapunov exponent via the QR decomposition method of Benettin et al.~\cite{Benettin1980}, integrating the variational equation for $T = 500$ time units after discarding a transient of $T_\text{trans} = 50$ time units, with QR renormalization every $\Delta t = 1.0$.
This yields $\lambda_{\max} \approx 0.3$, corresponding to a Lyapunov time of $T_\lambda = 1/\lambda_{\max} \approx 3.3$ time units.

This implies that even a perfect model will exhibit $\mathcal{O}(1)$ trajectory deviations from the truth on timescales of a few Lyapunov times when initialized from nearby (but not identical) states.
Consequently, we argue that trajectory-level metrics are informative only within ${\sim}1$--$2\,T_\lambda$ and that equation-level metrics---support recovery and coefficient error---are the appropriate criteria for evaluating discovered models in this regime.


%% ============================================================
\section{Methods}
\label{sec:methods}

\subsection{Data Generation}
\label{sec:data}

We generate training data using a ``short burst'' protocol: $n_\text{traj} = 250$ independent trajectories of length $T = 5.0$ time units, each integrated with a fourth/fifth-order Runge--Kutta method (SciPy \texttt{solve\_ivp} with \texttt{rtol}=$10^{-9}$, \texttt{atol}=$10^{-12}$) at a time step of $\Delta t = 0.01$.
Initial conditions are drawn uniformly from $[0.5,\, 2.0]^4$ for each state variable, providing broad coverage of the attractor basin.

The short-burst protocol avoids the need for a single long trajectory (which would sample the attractor statistically but provide less independent information per unit of data) and is more representative of experimental settings where repeated short acquisitions are available.

\subsection{Noise Model}
\label{sec:noise}

We corrupt the state measurements with additive Gaussian noise:
\begin{equation}
\tilde{\bX} = \bX + \eta\,\boldsymbol{\sigma}_X \odot \boldsymbol{\varepsilon}, \quad \boldsymbol{\varepsilon} \sim \mathcal{N}(0, \mathbf{I}),
\label{eq:noise}
\end{equation}
where $\eta \in \{10^{-4}, 10^{-3}, 10^{-2}, 10^{-1}, 1, 10\}$ is the noise amplitude, $\boldsymbol{\sigma}_X$ is the empirical standard deviation of the clean data computed per state variable, and $\odot$ denotes elementwise multiplication.
This normalization ensures that $\eta$ represents the noise-to-signal ratio uniformly across all four state variables.

\subsection{Candidate Libraries}
\label{sec:libraries}

We consider three candidate libraries of increasing complexity.
The physics-informed library encodes substantial domain knowledge and should be understood as a \emph{strong prior}---it is not representative of a blind discovery setting, but rather of the situation where a domain scientist has narrowed the functional forms and wishes to determine the precise sparse coefficients.

\paragraph{Physics-informed library ($p = 10$ features, strong prior).}
This library includes the constant function $1$, the four state variables $P$, $S$, $Z$, the trigonometric functions $\cos C$ and $\sin C$, the three bilinear-trigonometric terms $ZS\cos C$, $ZP\cos C$, $PS\cos C$, and the rational-trigonometric term $(PZ/S)\sin C$.
The ten functions span the CLW dynamics exactly, meaning the ground truth lies in the column space of the library.
Including this term requires prior knowledge of the rational coupling structure.

\paragraph{Extended library ($p = 46$ features, weak prior).}
Starting from the same ten basis functions, we augment the library with all pairwise cross-products of the nine non-constant basis functions, yielding $10 + \binom{9}{2} = 10 + 36 = 46$ features.
This more closely approximates a scenario with weaker prior knowledge, where a practitioner includes ``all plausible degree-2 terms'' and relies on sparsity to select the correct subset.

\paragraph{Incomplete library ($p = 9$ features, misspecified prior).}
We remove the rational term $(PZ/S)\sin C$ from the physics-informed library, leaving nine features.
Because the true $\dot{C}$ equation contains this term, the dynamics are no longer representable.
This ablation quantifies the penalty of library misspecification.

\subsection{The SINDy Algorithm}
\label{sec:sindy}

We use the Sequential Thresholded Least Squares (STLSQ) optimizer~\cite{Brunton2016} as implemented in PySINDy~\cite{deSilva2020,Kaptanoglu2022}.
STLSQ alternates between least-squares fitting and hard thresholding of small coefficients below a threshold $\lambda$.
We set the ridge regularization parameter $\alpha = 0$ and do not normalize library columns.

\paragraph{Model selection.}
For each noise level $\eta$ and seed, we sweep the STLSQ threshold over 25 logarithmically spaced values from $10^{-6}$ to $10^{0}$.
The best model is selected by minimizing a penalized score:
\begin{equation}
\text{score} = \log(\text{MSE}) + w_\text{nnz} \cdot \text{nnz},
\label{eq:model_selection}
\end{equation}
where MSE is the mean squared prediction error on the training data, nnz is the number of nonzero coefficients, and $w_\text{nnz} = 2 \times 10^{-3}$ is the sparsity penalty weight.
The logarithmic MSE ensures that the score favors order-of-magnitude improvements in fit quality over incremental gains from additional terms.
The weight $w_\text{nnz} = 2\times 10^{-3}$ was chosen so that a single additional term must reduce MSE by a factor of $e^{-0.002} \approx 0.998$ to be retained---a deliberately permissive threshold that slightly favors inclusion over exclusion.
We verify in Section~\ref{sec:threshold_path} that this criterion selects models near the ``knee'' of the threshold path (where further sparsification abruptly increases MSE) and that replacing training MSE with a held-out validation MSE yields indistinguishable selected models.
The constant term ``$1$'' is constrained to appear only in the $\dot{C}$ equation, reflecting the known structure.

\subsection{Derivative Estimation Regimes}
\label{sec:derivatives}

We compare three derivative estimation strategies of decreasing idealization:

\begin{enumerate}
\item \textbf{Analytic RHS (``oracle'')}: The known right-hand side of the CLW system~\eqref{eq:clw_P}--\eqref{eq:clw_C} is evaluated at the (possibly noisy) state measurements $\tilde{\bX}$ to obtain $\dot{\tilde{\bX}} = f(\tilde{\bX})$.
This provides exact derivatives \emph{only when the states are clean}; under noise, the derivatives inherit state-measurement error through the nonlinear RHS.
We stress that this regime does not provide noise-free derivatives when $\eta > 0$.

\item \textbf{Numerical finite differences (FD)}: Derivatives are computed from the noisy state measurements using second-order central differences via \texttt{numpy.gradient}.
No smoothing is applied.

\item \textbf{Savitzky--Golay smoothed finite differences (SFD)}: PySINDy's \texttt{SmoothedFiniteDifference} method is used, which applies a Savitzky--Golay filter~\cite{Savitzky1964} to smooth the state data before differentiating.
Unless otherwise noted, the default configuration uses a window length of 21, polynomial order 3, and differentiation order 2.
In this regime, SINDy handles differentiation internally on a per-trajectory basis.
\end{enumerate}

\subsection{Evaluation Metrics}
\label{sec:metrics}

We evaluate discovered models using the following metrics:

\paragraph{Number of nonzero coefficients (nnz).}
The total count of nonzero entries in $\hat{\bXi}$.
The true model has $k^\star = 8$.

\paragraph{Coefficient relative $L_2$ error.}
\begin{equation}
\epsilon_\text{coef} = \frac{\norm{\hat{\bXi} - \bXi^\star}_F}{\norm{\bXi^\star}_F},
\end{equation}
where $\bXi^\star$ is the ground-truth coefficient matrix aligned to the candidate library, and $\norm{\cdot}_F$ denotes the Frobenius norm.

\paragraph{True positive rate (TPR).}
The fraction of truly nonzero coefficients that are correctly identified as nonzero: $\text{TPR} = |\hat{S} \cap S^\star| / |S^\star|$, where $S^\star$ and $\hat{S}$ are the true and estimated support sets.

\paragraph{False positive rate (FPR).}
The fraction of truly zero coefficients that are incorrectly identified as nonzero: $\text{FPR} = |\hat{S} \setminus S^\star| / |S^0|$, where $S^0$ is the set of truly zero entries.

\paragraph{Exact support match.}
A binary indicator: $\hat{S} = S^\star$.

\paragraph{Vector field error.}
The relative $L_2$ error of the model's predicted right-hand side evaluated on the test data:
\begin{equation}
\epsilon_\text{vf} = \frac{\norm{\hat{f}(\bX_\text{test}) - f^\star(\bX_\text{test})}_F}{\norm{f^\star(\bX_\text{test})}_F}.
\end{equation}

All metrics are reported as means $\pm$ standard deviations over five independent random seeds.


%% ============================================================
\section{Experimental Design}
\label{sec:experiments}

We organize the study around six experiment groups, chosen to support a compact benchmark narrative rather than an exhaustive methods paper.

\paragraph{Experiment 1: Derivative regime comparison.}
For each of the three derivative estimation methods (analytic RHS, FD, SFD) and each noise level $\eta \in \{10^{-4}, 10^{-3}, 10^{-2}, 10^{-1}, 1, 10\}$, we fit a SINDy model using the physics-informed library and evaluate all metrics.
This is the primary experiment and isolates the effect of derivative quality on recovery.

\paragraph{Experiment 2: Library ablation.}
We compare the physics-informed, extended, and incomplete libraries.
In the noiseless case with analytic RHS on clean states, both the physics-informed and extended libraries achieve exact recovery (Table~\ref{tab:extended_clean}).
Under noise, the extended library is tested with analytic RHS derivatives across the full $\eta$ sweep to assess how library size interacts with noise (Table~\ref{tab:consolidated}).
The incomplete library is tested similarly to quantify misspecification effects.
For the extended library under noise, we additionally report which specific spurious terms are most frequently selected (Section~\ref{sec:false_positives}).

\paragraph{Experiment 3: Model-selection diagnostics.}
\label{sec:exp_diagnostics}
We examine the threshold-path behavior at $\eta = 0.01$ (physics-informed library, analytic RHS) to visualize how nnz, coefficient error, and the model-selection score vary across all 25 STLSQ thresholds.
We also compare train-MSE-based model selection with held-out validation-MSE-based selection across all seeds and noise levels (Section~\ref{sec:threshold_path}).

\paragraph{Experiment 4: Sample-size ablation.}
At a fixed moderate noise level $\eta = 0.01$ with analytic RHS derivatives and the physics-informed library ($p=10$), we sweep the number of training trajectories $n_\text{traj} \in \{25, 50, 100, 250, 500\}$ to assess data efficiency.
We pair this with a BIC-based variant to test whether the observed trend is an artifact of the fixed sparsity penalty.

\paragraph{Experiment 5: SFD sensitivity.}
We sweep the SFD window length $w \in \{5, 11, 21, 41\}$ and polynomial order $q \in \{2, 3, 5\}$ at three noise levels $\eta \in \{10^{-3}, 10^{-2}, 10^{-1}\}$ to characterize sensitivity to smoothing hyperparameters.

\paragraph{Experiment 6: Chaos-aware trajectory evaluation.}
We use short-horizon rollouts and Lyapunov-time normalization to illustrate why long-horizon trajectory agreement is not a reliable primary metric on CLW, even for the true model.


%% ============================================================
\section{Results}
\label{sec:results}

\subsection{Exact Recovery in the Noiseless Regime}

With clean data and analytic RHS derivatives (i.e.\ exact derivatives on exact states), both the physics-informed (10 features) and extended (46 features) libraries achieve exact sparse coefficient recovery (Table~\ref{tab:extended_clean}).
The discovered coefficient matrix matches the ground truth to machine precision ($\epsilon_\text{coef} < 4 \times 10^{-16}$), with $\text{nnz} = 8 = k^\star$, $\text{TPR} = 1.0$, $\text{FPR} = 0.0$, and exact support match in all five seeds.
This confirms that STLSQ can identify the correct sparse structure from the candidate library when given noiseless data and perfect derivatives, even when the library is substantially overspecified ($p=46$ vs.\ $k^\star=8$).
We note that this result depends on the physics-informed library being a \emph{strong prior}: the 10-feature library contains exactly the true terms, while even the extended library is constructed by augmenting those same physics-informed basis functions rather than by generic polynomial expansion.

\begin{table}[t]
\centering
\caption{Clean recovery with analytic RHS derivatives on noise-free states (mean $\pm$ std, 5 seeds). Both libraries achieve exact coefficient recovery.}
\label{tab:extended_clean}
\begin{tabular}{lcccccc}
\toprule
Library & nnz & $\epsilon_\text{coef}$ & TPR & FPR & Exact & $\epsilon_\text{vf}$ \\
\midrule
Physics-informed ($p{=}10$) & 8 $\pm$ 0 & $3.76 \times 10^{-16}$ & 1.0 & 0.0 & 5/5 & $4.08 \times 10^{-16}$ \\
Extended ($p{=}46$) & 8 $\pm$ 0 & $3.76 \times 10^{-16}$ & 1.0 & 0.0 & 5/5 & $4.08 \times 10^{-16}$ \\
\bottomrule
\end{tabular}
\end{table}


\subsection{Derivative Regime Comparison Under Noise}

Table~\ref{tab:consolidated} presents the central result: a comparison of the three derivative estimation regimes across noise levels.
Figures~\ref{fig:rel_l2} and~\ref{fig:tpr} visualize the coefficient error and TPR trends.

\begin{table}[t]
\centering
\caption{Consolidated derivative regime comparison (mean $\pm$ std, 5 seeds). Physics-informed library ($p{=}10$).}
\label{tab:consolidated}
\small
\begin{tabular}{llcccccc}
\toprule
Regime & $\eta$ & nnz & $\epsilon_\text{coef}$ & TPR & FPR & $\epsilon_\text{vf}$ \\
\midrule
\multirow{6}{*}{\shortstack[l]{Analytic\\RHS}}
& $10^{-4}$ & 10.6 & 0.187 & 1.00 & 0.08 & 0.363 \\
& $10^{-3}$ & 10.2 & 0.469 & 1.00 & 0.07 & 0.843 \\
& $10^{-2}$ & 12.0 & 0.395 & 0.93 & 0.14 & 0.362 \\
& $10^{-1}$ & 12.2 & 0.470 & 0.88 & 0.16 & 0.448 \\
& $1$       & 4.0  & 1.007 & 0.13 & 0.09 & 0.962 \\
& $10$      & 0.8  & 0.935 & 0.10 & 0.00 & 0.990 \\
\midrule
\multirow{6}{*}{FD}
& $10^{-4}$ & 12.0 & 0.117 & 1.00 & 0.13 & 0.203 \\
& $10^{-3}$ & 10.4 & 0.027 & 1.00 & 0.08 & 0.030 \\
& $10^{-2}$ & 11.2 & 0.346 & 0.98 & 0.11 & 0.350 \\
& $10^{-1}$ & 6.2  & 0.625 & 0.58 & 0.05 & 0.555 \\
& $1$       & 2.8  & 1.079 & 0.08 & 0.07 & 0.988 \\
& $10$      & 6.8  & 3.665 & 0.03 & 0.21 & 1.842 \\
\midrule
\multirow{6}{*}{SFD}
& $10^{-4}$ & 11.0 & 0.026 & 1.00 & 0.09 & 0.024 \\
& $10^{-3}$ & 11.0 & 0.027 & 1.00 & 0.09 & 0.025 \\
& $10^{-2}$ & 11.4 & 0.112 & 1.00 & 0.11 & 0.084 \\
& $10^{-1}$ & 2.6  & 0.711 & 0.33 & 0.00 & 0.685 \\
& $1$       & 3.8  & 0.942 & 0.35 & 0.03 & 0.779 \\
& $10$      & 2.4  & 1.436 & 0.08 & 0.06 & 1.058 \\
\bottomrule
\end{tabular}
\end{table}

\begin{figure}[t]
\centering
\includegraphics[width=0.85\textwidth]{fig2_rel_l2_vs_eta.pdf}
\caption{Coefficient relative $L_2$ error versus noise level $\eta$ for the three derivative estimation regimes. SFD provides the best low-noise recovery but degrades sharply at $\eta \geq 0.1$.}
\label{fig:rel_l2}
\end{figure}

\begin{figure}[t]
\centering
\includegraphics[width=0.85\textwidth]{fig3_tpr_vs_eta.pdf}
\caption{True positive rate versus noise level $\eta$. All regimes maintain $\text{TPR} = 1.0$ for $\eta \leq 0.01$; SFD degrades most rapidly at higher noise.}
\label{fig:tpr}
\end{figure}

Several patterns are worth emphasizing.

\paragraph{Low noise ($\eta \leq 10^{-2}$).}
All three regimes retain most or all of the true support, but they are not equally accurate.
SFD is best in this regime, with $\epsilon_\text{coef} = 0.026$--$0.112$ and perfect mean TPR through $\eta = 10^{-2}$.
Raw finite differences are competitive at $\eta = 10^{-3}$.
The analytic RHS regime is noticeably worse than either data-driven alternative at low noise because evaluating $f(\tilde{\bX})$ on noisy states pushes the state perturbation through products, trigonometric terms, and the rational factor $PZ/S$.
In other words, the ``oracle'' regime is only idealized when the states are clean; once the states are noisy, it is a nonlinear error-propagation experiment rather than a gold standard.

\paragraph{Moderate noise ($\eta = 10^{-1}$).}
This is the practical cliff for the present benchmark.
SFD collapses to overly sparse models ($\text{nnz}=2.6$, mean TPR $0.33$), while the analytic RHS and FD regimes still retain some structure but no longer recover the equations cleanly.
The exact failure point is system-specific, but the qualitative message is robust: for CLW, standard derivative handling does not buy much usable headroom beyond $\eta \approx 10^{-2}$.

\paragraph{High noise ($\eta \geq 1$).}
Recovery fails across all regimes.
The analytic RHS and SFD settings collapse toward trivial models, while FD becomes numerically unstable enough that the coefficient and vector-field errors exceed the zero-model baseline.
This is the right place to stop claiming recovery and start treating the output as a failure mode.


\subsection{Library Ablation}

\paragraph{Extended library under noise.}
Table~\ref{tab:library_ablation} compares the physics-informed (analytic RHS) and extended (analytic RHS) libraries under noise.
In the clean case, both achieve exact recovery.
Under noise, the extended library does not destroy identifiability outright, but it creates more room for unstable false positives.
That is the more realistic lesson for practice: a moderately overcomplete library may still work in easy regimes, yet the price is paid in interpretability once noise enters.

\subsubsection{False-Positive Term Identification}
\label{sec:false_positives}

To understand \emph{which} spurious terms the extended library selects under noise, we record the identity of every false-positive term across all 5 seeds and 5 noise levels ($\eta\in\{10^{-4},...,1\}$), yielding 25 total runs.
Table~\ref{tab:false_positives} reports the most frequently selected non-true terms.
The dominant false positives are concentrated in the $\dot{C}$ equation, where the linear terms $S$~(84\%), $P$~(80\%), and $\sin C$~(68\%) are most frequently included spuriously.
This pattern is physically interpretable: the true $\dot{C}$ equation contains the rational-trigonometric term $(PZ/S)\sin C$, which is strongly correlated with its constituent sub-products $P$, $S$, $Z$, $\sin C$ when evaluated on the chaotic attractor.
Under noise, regression variance ``leaks'' into these correlated features.
The $\dot{P}$ and $\dot{S}$ equations show much lower false-positive rates ($\leq$20\%), consistent with their simpler nonlinear structure.
No false-positive term appears in more than 84\% of runs, indicating that individual spurious inclusions are somewhat unstable across seeds and noise levels.

\begin{table}[t]
\centering
\caption{Most frequently selected false-positive terms in the extended library ($p=46$), across 5 seeds $\times$ 5 noise levels (analytic RHS). Terms appearing in fewer than 2/25 runs are omitted.}
\label{tab:false_positives}
\small
\begin{tabular}{llcc}
\toprule
Equation & Spurious term & Count (/25) & Fraction \\
\midrule
$\dot{C}$ & $S$       & 21 & 0.84 \\
$\dot{C}$ & $P$       & 20 & 0.80 \\
$\dot{C}$ & $\sin C$  & 17 & 0.68 \\
$\dot{C}$ & $\cos C$  &  7 & 0.28 \\
$\dot{P}$ & $S$       &  5 & 0.20 \\
$\dot{P}$ & $\cos C$  &  4 & 0.16 \\
$\dot{S}$ & $P$       &  4 & 0.16 \\
$\dot{P}$ & $\sin C$  &  2 & 0.08 \\
$\dot{C}$ & $Z$       &  2 & 0.08 \\
\bottomrule
\end{tabular}
\end{table}

\paragraph{Incomplete library.}
The incomplete library, missing $(PZ/S)\sin C$, produces a persistent coefficient-error floor near $0.48$ (Table~\ref{tab:library_ablation}).
This is one of the clearest findings in the paper.
Even very low measurement noise cannot compensate for a missing structural term.
The optimizer simply redistributes the missing dynamics onto available features, keeping TPR high for the recoverable terms while preserving a large irreducible model error.
For this benchmark, library misspecification is therefore more damaging than moving from clean derivatives to mildly noisy derivatives.

\begin{table}[t]
\centering
\caption{Library ablation: physics-informed vs.\ extended vs.\ incomplete (analytic RHS derivatives, mean over 5 seeds).}
\label{tab:library_ablation}
\small
\begin{tabular}{llccccc}
\toprule
Library & $\eta$ & nnz & $\epsilon_\text{coef}$ & TPR & FPR & $\epsilon_\text{vf}$ \\
\midrule
\multirow{3}{*}{Physics ($p{=}10$)}
& $10^{-4}$ & 10.6 & 0.187 & 1.00 & 0.08 & 0.363 \\
& $10^{-2}$ & 12.0 & 0.395 & 0.93 & 0.14 & 0.362 \\
& $10^{-1}$ & 12.2 & 0.470 & 0.88 & 0.16 & 0.448 \\
\midrule
\multirow{3}{*}{Extended ($p{=}46$)}
& $10^{-4}$ & 10.0 & 0.495 & 1.00 & 0.06 & 0.960 \\
& $10^{-2}$ & 11.6 & 0.291 & 0.98 & 0.12 & 0.305 \\
& $10^{-1}$ & 12.0 & 0.468 & 0.88 & 0.16 & 0.447 \\
\midrule
\multirow{3}{*}{Incomplete ($p{=}9$)}
& $10^{-4}$ & 12.4 & 0.480 & 1.00 & 0.19 & 0.444 \\
& $10^{-2}$ & 11.6 & 0.493 & 0.94 & 0.17 & 0.444 \\
& $10^{-1}$ & 12.0 & 0.421 & 1.00 & 0.17 & 0.447 \\
\bottomrule
\end{tabular}
\end{table}

\begin{figure}[t]
\centering
\includegraphics[width=0.85\textwidth]{fig5_library_ablation.pdf}
\caption{Library ablation: coefficient error across noise levels for the three library types. The incomplete library has a persistent error floor at ${\sim}0.48$ due to the missing rational term.}
\label{fig:library_ablation}
\end{figure}


\subsection{Sample-Size Ablation}

Table~\ref{tab:sample_size} reports the effect of training set size at fixed $\eta = 0.01$ with analytic RHS derivatives and the physics-informed library ($p=10$).
All configurations maintain high TPR, but the coefficient error and FPR vary non-monotonically.
At $n_\text{traj} = 25$, the best mean results are obtained ($\epsilon_\text{coef} = 0.107$, exact support in 1/5 seeds), while larger datasets show higher mean coefficient errors.

This counterintuitive pattern does not appear to be a small-sample artifact.
Instead, it arises from the interaction between model selection and correlated features on the attractor.
The penalized score~\eqref{eq:model_selection} selects the threshold $\lambda$ that minimizes $\log(\text{MSE}) + w_\text{nnz}\cdot\text{nnz}$.
With more training data, weakly correlated but physically spurious features become easier to justify by pure prediction error, so the selected model can become less sparse even though the data are more informative.
This is exactly the kind of result that a benchmark paper should report plainly: more data do not automatically improve equation recovery when the library contains correlated alternatives.

\paragraph{BIC model selection does not resolve the anomaly.}
To test whether the non-monotonic behavior is an artifact of the fixed sparsity penalty, we re-run the same sample-size sweep using the Bayesian Information Criterion: $\text{BIC} = n\,\log(\text{MSE}) + k\,\log n$, where $n$ is the total number of data points and $k = \text{nnz}$.
Table~\ref{tab:sample_size_bic} shows that the anomaly persists and is in fact more pronounced under BIC:
the FPR rises from 0.11 at $n_\text{traj}=25$ to 0.23 at $n_\text{traj}=500$, and no seed achieves exact support for $n_\text{traj}\geq 50$.

This suggests that the problem is structural rather than merely heuristic.
The spurious terms are genuinely predictive on the attractor, so a generic information criterion has little reason to reject them.
For CLW, sparsity alone is not a sufficient proxy for physical correctness once the library contains correlated surrogates for the true rational interaction.

\begin{table}[t]
\centering
\caption{BIC model selection: sample-size ablation ($\eta = 0.01$, analytic RHS, physics-informed library $p=10$, 5 seeds). Compare with the fixed-penalty results in Table~\ref{tab:sample_size}.}
\label{tab:sample_size_bic}
\begin{tabular}{rccccc}
\toprule
$n_\text{traj}$ & nnz & $\epsilon_\text{coef}$ & TPR & FPR & Exact \\
\midrule
25  & $11.6 \pm 2.6$ & $0.108 \pm 0.163$ & 1.00 & 0.11 & 1/5 \\
50  & $14.2 \pm 2.6$ & $0.241 \pm 0.202$ & 1.00 & 0.19 & 0/5 \\
100 & $14.6 \pm 1.5$ & $0.235 \pm 0.204$ & 1.00 & 0.21 & 0/5 \\
250 & $14.8 \pm 2.4$ & $0.339 \pm 0.170$ & 0.98 & 0.22 & 0/5 \\
500 & $15.0 \pm 2.0$ & $0.383 \pm 0.151$ & 0.98 & 0.23 & 0/5 \\
\bottomrule
\end{tabular}
\end{table}

\begin{table}[t]
\centering
\caption{Sample-size ablation ($\eta = 0.01$, analytic RHS, physics-informed library $p=10$, 5 seeds).}
\label{tab:sample_size}
\begin{tabular}{rccccc}
\toprule
$n_\text{traj}$ & nnz & $\epsilon_\text{coef}$ & TPR & FPR & Exact \\
\midrule
25  & $11.2 \pm 2.3$ & $0.107 \pm 0.162$ & 1.00 & 0.10 & 1/5 \\
50  & $12.2 \pm 2.6$ & $0.235 \pm 0.202$ & 1.00 & 0.13 & 1/5 \\
100 & $11.8 \pm 2.6$ & $0.231 \pm 0.204$ & 1.00 & 0.12 & 1/5 \\
250 & $11.8 \pm 1.3$ & $0.330 \pm 0.169$ & 0.98 & 0.13 & 0/5 \\
500 & $12.6 \pm 2.1$ & $0.376 \pm 0.150$ & 0.98 & 0.15 & 0/5 \\
\bottomrule
\end{tabular}
\end{table}


\subsection{SFD Sensitivity}

Figure~\ref{fig:sfd} shows the SFD hyperparameter sensitivity.
The key findings are:

\begin{itemize}
\item At low noise ($\eta = 10^{-3}$), short windows ($w = 5$) perform best, achieving the lowest $\epsilon_\text{coef}$ (${\sim}0.011$) and exact support match in 2/5 seeds. Excessively long windows ($w = 41$) introduce smoothing bias, inflating nnz to ${\sim}20$ and $\epsilon_\text{coef}$ to ${\sim}0.11$--$0.14$.
\item At moderate noise ($\eta = 10^{-2}$), intermediate windows ($w = 11$--$21$) provide the best balance, with $\epsilon_\text{coef} \approx 0.02$--$0.05$ and $\text{TPR} = 1.0$.
\item At high noise ($\eta = 10^{-1}$), no SFD configuration recovers more than $\text{TPR} = 0.88$, and coefficient error exceeds $0.48$ across all window/order combinations. The Savitzky--Golay filter cannot rescue recovery at this noise level.
\item Higher polynomial orders ($q = 5$) are generally beneficial at larger windows, mitigating the smoothing bias that afflicts low-order fits.
\end{itemize}

\begin{figure}[t]
\centering
\includegraphics[width=0.85\textwidth]{figS1_sfd_sensitivity.pdf}
\caption{SFD sensitivity: coefficient error as a function of window length and polynomial order, at three noise levels.}
\label{fig:sfd}
\end{figure}


\subsection{Model-Selection Diagnostics}
\label{sec:threshold_path}

To justify the model-selection criterion~\eqref{eq:model_selection}, we examine the threshold path: the trajectory of key metrics across all 25 STLSQ thresholds at $\eta=0.01$ (physics-informed library, analytic RHS, seed~0).
Figure~\ref{fig:threshold_path} shows that as the threshold increases from $10^{-6}$ to $10^{0}$, nnz decreases monotonically (from 37 to 4) while MSE increases; the penalized score exhibits a clear minimum near the ``knee'' of this trade-off.
The selected model at this minimum (threshold $\lambda\approx 0.056$) has $\text{nnz}=11$ (8 true $+$ 3 spurious) with $\epsilon_\text{coef}\approx 0.10$.
Notably, the exact-support model ($\text{nnz}=8$, $\lambda\approx 0.56$) achieves a lower coefficient error ($\epsilon_\text{coef}\approx 0.05$) but a slightly worse score because the MSE jump from removing 3 terms exceeds their sparsity penalty---illustrating the known trade-off of penalized model selection.

We also compare train-MSE-based selection with held-out validation on a 20\% trajectory split.
Across all 5 seeds $\times$ 5 noise levels, the mean coefficient error differs by less than 10\% between the two selection strategies ($\bar{\epsilon}_\text{train}=0.51$ vs.\ $\bar{\epsilon}_\text{val}=0.50$), and the mean nnz is comparable ($\bar{\text{nnz}}_\text{train}=9.8$ vs.\ $\bar{\text{nnz}}_\text{val}=10.1$).
The two methods select the same STLSQ threshold in the majority of seed/noise combinations.
This confirms that the training-MSE-based score is not overfitting the threshold choice for this data volume ($n_\text{traj}=250$, 500 time steps each).

\begin{figure}[t]
\centering
\includegraphics[width=0.85\textwidth]{fig_threshold_path.pdf}
\caption{Threshold-path diagnostics ($\eta=0.01$, physics-informed library, analytic RHS, seed~0). Left: nnz vs.\ threshold. Center: coefficient error vs.\ threshold. Right: penalized score (train and validation) vs.\ threshold. The vertical dashed line marks the selected threshold.}
\label{fig:threshold_path}
\end{figure}


\subsection{The Chaos Evaluation Problem}

Figure~\ref{fig:chaos} illustrates the central evaluation issue for this benchmark: two trajectories of the true CLW system initialized with $\Delta C = 10^{-6}$ diverge to $\mathcal{O}(1)$ differences within about $3\,T_\lambda$.
That divergence is not a failure of identification; it is a property of the underlying dynamics.
Long-horizon rollout mismatch therefore cannot be the main criterion for acceptance or rejection of a discovered model here.

Figure~\ref{fig:error_vs_lyapunov} makes the supporting point: in the low-noise SFD regime, the learned model follows the reference dynamics for about one to two Lyapunov times before the expected divergence, whereas noisier models separate almost immediately.
Short-horizon agreement is therefore useful as a sanity check, but support recovery and coefficient error remain the better primary metrics.

\begin{figure}[t]
\centering
\includegraphics[width=0.85\textwidth]{figS2_chaos_demo.pdf}
\caption{Chaos demonstration: two CLW trajectories with $\Delta C = 10^{-6}$ in initial conditions. The \emph{true model} produces $\mathcal{O}(1)$ deviations within ${\sim}3\,T_\lambda$, exposing the inadequacy of long-horizon trajectory error as a model quality metric.}
\label{fig:chaos}
\end{figure}

\begin{figure}[t]
\centering
\includegraphics[width=0.85\textwidth]{fig4_error_vs_lyapunov_time.pdf}
\caption{Trajectory error vs.\ time (in Lyapunov units $T_\lambda \approx 3.3$) for discovered models in the three derivative regimes at selected noise levels. Low-noise SFD models track the truth for ${\sim}1$--$2\,T_\lambda$ before the inevitable chaotic divergence.}
\label{fig:error_vs_lyapunov}
\end{figure}

\begin{figure}[t]
\centering
\includegraphics[width=\textwidth]{fig1_clw_overview.pdf}
\caption{CLW system overview. (Left)~Time series of the four state variables $P$, $S$, $Z$, $C$. (Right)~Three-dimensional phase portrait in $(P, S, Z)$ space, colored by time, showing the chaotic attractor.}
\label{fig:clw_overview}
\end{figure}


%% ============================================================
\section{Discussion}
\label{sec:discussion}

\subsection{What This Benchmark Shows}

The paper supports four main claims.
First, baseline STLSQ-SINDy can recover the CLW equations exactly in the idealized regime where both the states and derivatives are effectively exact and the library contains the true functional form.
Second, for this system the practically useful noise range is limited: coefficient recovery remains credible through about $\eta = 10^{-2}$ and breaks down near $\eta = 10^{-1}$ regardless of derivative strategy.
Third, library choice matters at least as much as derivative estimation.
An overcomplete library mainly harms interpretability through false positives, while an incomplete library can impose a large error floor even when the data are almost clean.
Fourth, chaotic rollout divergence should be interpreted carefully; it is expected behavior on CLW and does not by itself refute a good local equation model.

These claims are narrower than the claims typically made in methodological machine-learning papers, but that narrower scope is a strength here.
The value of the study lies in calibration: a reader can see what a standard pipeline does on a nontrivial plasma-motivated system before trying more elaborate identification methods.

\subsection{Practical Implications}

For practitioners, the results suggest a simple ordering of priorities.
If a plausible physics-informed library is available, getting the library structure right is more important than marginal improvements in threshold tuning.
If the measurements are noisy, smoothed finite differences are worthwhile in the low-noise regime, but they should not be expected to rescue recovery once the derivative signal is badly contaminated.
And if the goal is scientific interpretation rather than forecasting, support recovery and coefficient error should be treated as the headline metrics.

The sample-size and BIC ablations carry a less intuitive lesson.
More data do not necessarily make the selected equation cleaner when the library contains correlated surrogate terms.
On CLW, additional data can strengthen the evidence for terms that are predictive on the attractor but still physically wrong.
That is a caution against assuming that generic information criteria are enough to separate physical structure from statistical convenience.

\subsection{Limitations}

We acknowledge several limitations of this study:

\begin{itemize}
\item \textbf{Single system, single parameter set}: Our experiments are restricted to the CLW model at $\Gamma_d=2$, $\gamma_z=0.8$, $d=2$. While this provides a complete controlled study, the quantitative noise thresholds and error magnitudes may not generalize to other chaotic ODE systems or other CLW parameter regimes.

\item \textbf{Reduced model vs.\ real plasma turbulence}: The CLW system is a low-dimensional ($n=4$) reduction of tokamak turbulence dynamics. Real plasma dynamics involve spatial degrees of freedom ($n\gg 4$), multiple time scales, closures with uncertain coefficients, and diagnostic noise that is non-Gaussian, state-dependent, and temporally correlated. Our additive Gaussian noise model is idealized.

\item \textbf{Strong prior in library design}: The physics-informed library was designed with knowledge of the true dynamics, including the rational nonlinearity $PZ/S$. In practice, a domain expert may not know the relevant nonlinearities \emph{a priori}. The extended-library results partially address this, but even the extended library is built from physics-informed basis functions.

\item \textbf{Model selection remains heuristic}: The penalized score~\eqref{eq:model_selection} behaves reasonably and agrees well with held-out validation, but the sample-size ablation shows that no simple generic criterion resolves the correlated-feature problem on this attractor. A better selection rule may require additional physics priors or structured regularization.

\item \textbf{Baseline algorithm only}: We deliberately study a standard STLSQ pipeline. More advanced formulations---for example weak-form or Bayesian approaches~\cite{Messenger2021}---may improve robustness. The purpose of the present paper is to supply a transparent baseline against which such methods can be judged.
\end{itemize}


%% ============================================================
\section{Conclusion}
\label{sec:conclusion}

We have presented a controlled benchmark study of sparse equation discovery on the Chen--Lin--White reduced tokamak turbulence model.
The main outcome is not a new algorithmic claim, but a calibrated picture of what baseline STLSQ-SINDy can and cannot do on this chaotic system.

\begin{enumerate}
\item \textbf{Idealized recovery is possible}: with clean states, exact derivatives, and a library that contains the true functional forms, CLW is recovered exactly.

\item \textbf{The useful noise range is limited}: for this benchmark, recovery is credible up to about $\eta = 10^{-2}$ and degrades sharply near $\eta = 10^{-1}$ across all derivative regimes.

\item \textbf{Library structure is decisive}: overspecification mainly introduces false positives, whereas misspecification creates a large irreducible error floor even at negligible noise.

\item \textbf{Generic model-selection criteria have limits}: threshold-path and validation diagnostics support the chosen score, but neither that score nor BIC removes spurious correlated terms once they become predictively useful on the attractor.

\item \textbf{Chaotic systems should be judged locally and structurally}: short-horizon agreement is informative, but support recovery, coefficient error, and vector-field error are the more reliable metrics.
\end{enumerate}

CLW is still a small and idealized problem, but it contains enough structure to expose real difficulties: nonlinear coupling, a rational interaction term, correlated surrogate features, and sensitive dependence on initial conditions.
That makes it a useful intermediate benchmark between textbook dynamical systems and full plasma turbulence models.
We expect the main value of this paper to be as a baseline reference for future studies that test stronger differentiation schemes, better-structured libraries, or more robust sparse regression methods on the same system.

\subsection*{Code and Data Availability}

All code, data, and scripts to reproduce every table, figure, and analysis in this paper are available at \url{https://github.com/lucas-perrier/sindy-clw}.
The repository includes fixed random seeds, pinned package versions (\texttt{requirements.txt}), and a single entry point (\texttt{experiments/run\_all.py}) that regenerates all results.


%% ============================================================
%% References
%% ============================================================

\begin{thebibliography}{99}

\bibitem{Benettin1980}
G.~Benettin, L.~Galgani, A.~Giorgilli, J.-M.~Strelcyn,
Lyapunov characteristic exponents for smooth dynamical systems and for {H}amiltonian systems; a method for computing all of them,
Meccanica 15 (1980) 9--20.

\bibitem{Brunton2016}
S.L.~Brunton, J.L.~Proctor, J.N.~Kutz,
Discovering governing equations from data by sparse identification of nonlinear dynamical systems,
Proc.\ Natl.\ Acad.\ Sci.\ USA 113 (2016) 3932--3937.

\bibitem{Champion2019}
K.~Champion, B.~Lusch, J.N.~Kutz, S.L.~Brunton,
Data-driven discovery of coordinates and governing equations,
Proc.\ Natl.\ Acad.\ Sci.\ USA 116 (2019) 22445--22451.

\bibitem{Chartrand2011}
R.~Chartrand,
Numerical differentiation of noisy, nonsmooth data,
ISRN Appl.\ Math.\ 2011 (2011) 164564.

\bibitem{Chen1984}
L.~Chen, Z.~Lin, R.B.~White,
On the resonant three-wave interaction in magnetized plasma,
Phys.\ Fluids 27 (1984) 1242--1248.

\bibitem{Dam2017}
M.~Dam, M.~Br{\o}ns, J.~Juul~Rasmussen, V.~Naulin, J.S.~Hesthaven,
Sparse identification of a predator-prey system from simulation data of a convection model,
SIAM J.\ Appl.\ Dyn.\ Syst.\ 16 (2017) 2334--2351.

\bibitem{deSilva2020}
B.~de~Silva, K.~Champion, M.~Quade, J.-C.~Loiseau, J.N.~Kutz, S.L.~Brunton,
{PySINDy}: A {P}ython package for the sparse identification of nonlinear dynamical systems from data,
J.\ Open Source Softw.\ 5 (2020) 2104.

\bibitem{Fasel2022}
U.~Fasel, J.N.~Kutz, B.W.~Brunton, S.L.~Brunton,
Ensemble-{SINDy}: Robust sparse model discovery in the low-data, high-noise limit, with active learning and control,
Proc.\ R.\ Soc.\ A 478 (2022) 20210904.

\bibitem{Horton1999}
W.~Horton,
Drift waves and transport,
Rev.\ Mod.\ Phys.\ 71 (1999) 735--778.

\bibitem{Kaptanoglu2020}
A.A.~Kaptanoglu, K.D.~Morgan, C.J.~Hansen, S.L.~Brunton,
Characterizing magnetized plasmas with dynamic mode decomposition,
Phys.\ Plasmas 27 (2020) 032108.

\bibitem{Kaptanoglu2022}
A.A.~Kaptanoglu, B.M.~de~Silva, U.~Fasel, K.~Kaheman, A.J.~Goldschmidt, J.~Callaham, C.B.~Delahunt, Z.G.~Nicolaou, K.~Champion, J.-C.~Loiseau, J.N.~Kutz, S.L.~Brunton,
{PySINDy}: A comprehensive {P}ython package for robust sparse system identification,
J.\ Open Source Softw.\ 7 (2022) 3994.

\bibitem{Mangan2016}
N.M.~Mangan, S.L.~Brunton, J.L.~Proctor, J.N.~Kutz,
Inferring biological networks by sparse identification of nonlinear dynamics,
IEEE Trans.\ Mol.\ Biol.\ Multi-Scale Commun.\ 2 (2016) 52--63.

\bibitem{Messenger2021}
D.A.~Messenger, D.M.~Bortz,
Weak {SINDy}: Galerkin-based data-driven model selection,
Multiscale Model.\ Simul.\ 19 (2021) 1474--1497.

\bibitem{Pathak2018}
J.~Pathak, B.~Hunt, M.~Girvan, Z.~Lu, E.~Ott,
Model-free prediction of large spatiotemporally chaotic systems from data: A reservoir computing approach,
Phys.\ Rev.\ Lett.\ 120 (2018) 024102.

\bibitem{Rudy2017}
S.H.~Rudy, S.L.~Brunton, J.L.~Proctor, J.N.~Kutz,
Data-driven discovery of partial differential equations,
Sci.\ Adv.\ 3 (2017) e1602614.

\bibitem{Savitzky1964}
A.~Savitzky, M.J.E.~Golay,
Smoothing and differentiation of data by simplified least squares procedures,
Anal.\ Chem.\ 36 (1964) 1627--1639.

\bibitem{Schaeffer2017}
H.~Schaeffer,
Learning partial differential equations via data discovery and sparse optimization,
Proc.\ R.\ Soc.\ A 473 (2017) 20160446.

\bibitem{Strogatz2015}
S.H.~Strogatz,
Nonlinear Dynamics and Chaos, 2nd ed., Westview Press, 2015.

\bibitem{Vlachas2018}
P.R.~Vlachas, W.~Byeon, Z.Y.~Wan, T.P.~Sapsis, P.~Koumoutsakos,
Data-driven forecasting of high-dimensional chaotic systems with long short-term memory networks,
Proc.\ R.\ Soc.\ A 474 (2018) 20170844.

\end{thebibliography}

\end{document}
